\chapter{绪论}

\section{课题的研究背景与意义}

随着人口老龄化不断加深，慢性病管理已成为公共健康领域的长期议题。在常见慢性病中，高血压、高血糖和高血脂患者数量大、管理周期长、复发风险高，患者往往不能只依赖少数几次门诊完成健康管理。与急性疾病不同，慢性病管理更强调日常监测和持续干预，患者在真实生活中会频繁遇到诸如指标是否异常、饮食如何调整、药物需不需要复查、症状是否需要就医等具体问题\cite{ZWSG202601012}。这类问题看似琐碎，却直接影响持续治疗情况和病情控制效果。

传统医疗服务以线下诊疗为主，患者通常只有感到身体不适或按固定的时间安排才会选择就医，患者难以获得持续、及时的健康指导。同时，由于医疗资源分布不均以及医生数量相对有限，患者在非就诊期间获取专业医疗建议的渠道较为有限。所以很多患者选择去互联网搜索相关健康信息，但互联网上信息来源复杂、质量参差不齐，用户很难从中获得可靠的医学建议。

近年来，大语言模型（Large Language Model，LLM）发展很快，在理解和生成文本方面都有明显提升，通过大语言模型构建智能医疗问答系统成为可能。基于LLM的问答系统能够理解用户提出的问题，并生成相对流畅的回答，在医疗咨询等场景中展现出一定应用潜力。但是将通用LLM直接应用在医学领域时，会面临知识更新滞后、事实错误（幻觉）以及缺乏依据等问题\cite{YXQB202601003}。

基于这一背景，本文引入检索增强生成（Retrieval-Augmented Generation，RAG）技术，将外部医学知识库与大语言模型结合起来\cite{Ji_2023}。其核心思路是在模型生成答案前先检索相关医学资料，根据资料组织回答，而不是单纯依靠模型的“记忆”作答。这样做的价值主要体现在三个方面：第一，可以降低医学问答中的幻觉风险；第二，可以让回答建立在可追溯的证据基础上；第三，更便于系统随知识更新而持续维护。因而，围绕“三高”慢病场景研究并实现一套基于RAG的医学问答系统，既具有现实应用意义，也具有一定的技术研究价值。

\section{国内外研究现状}

\subsection{互联网问诊的发展现状}

随着互联网技术和智能手机的普及，互联网问诊模式出现并迅速发展。根据中国互联网络信息中心第56次《中国互联网络发展状况统计报告》显示，截至2025年底，中国网民规模达11.25亿，互联网普及率突破80\%\cite{CNNICReport56_2025}。互联网问诊依托其便捷高效、实时互动、及良好的隐私保护特性等优势，迅速普及，吸引了大量互联网企业、传统医疗机构及新兴在线医疗平台积极参与\cite{WSGL201702004}。目前，国内的互联网医疗代表性平台有微医、好大夫在线、京东健康、平安健康和阿里健康等，其主要提供医生在线问诊、健康信息发布以及购药等服务。这些平台更多是作为传统线下医疗服务的补充形式，在一定程度上提升了患者获取医疗信息的便利性。不过，平台提供的信息质量参差不齐，不同来源之间可能存在内容差异，影响用户判断；在线问诊以人工问答为主，受限于有限的医疗资源，患者还是会面临回复不及时等问题。

国外的医疗平台则更注重于远程医疗。例如，瑞典的KRY和美国Teladoc Health、Doctor on Demand等平台实行线上预约并利用会议视频实时问诊，让患者在家里即可享受到专业的医疗问诊\cite{WRKX201206003}。但这类服务仍有很多不足，首先是患者可能需要支出额外的费用，以享受远程问诊的服务；而且平台签约医生数量相对固定，高峰时段还需预约排队；如果涉及身体检查的情况，仍需要去线下医院就诊。

现阶段互联网问诊仍主要依赖人工医生答复，其服务能力、响应效率和服务成本都受到现实医疗资源的制约。在这一背景下，如何借助自然语言处理与智能问答技术，为用户提供更及时、更稳定的医学信息支持，逐渐成为医学人工智能研究的重要问题，因此也推动了医学问答模型的持续发展。

\subsection{传统医学问答模型的发展现状}

早期医学自然语言处理模型主要集中在医学术语识别、实体抽取和文献理解等任务上，代表性工作包括BioBERT、BlueBERT、ClinicalBERT和PubMedBERT等。这类模型大多建立在BERT框架之上，通过PubMed文献、临床记录或生物医学语料继续预训练，使模型在医学词汇和专业语义建模方面明显优于通用语言模型。其主要价值在于提升医学文本表示能力，为后续问答、分类和关系抽取等任务提供基础，但这类模型主要是以理解任务为主。

随着大模型技术的发展，医学问答研究进一步转向生成式模型。例如，由Google提出的Med-PaLM和Med-PaLM 2模型，Med-PaLM系列模型是专门针对于医学领域的大语言模型，Med-PaLM通过医学领域微调和人工评估，使其在医学考试题、患者咨询和临床推理任务中相较于通用大语言模型有更好的表现，Med-PaLM也是第一个在美国医师执照考试样例问题中获得“合格”的模型\cite{Singhal_2025}。国内的科研机构也在积极推进大模型在医学领域的应用，也推出了许多针对于医学领域的大模型，例如：DoctorGLM、华佗GPT等。DoctorGLM是基于ChatGLM的中文问诊模型，其对通用大语言做了针对医疗领域进行微调与优化工作；香港中文大学和深圳大数据研究院的研究团队推出的华佗GPT-o1，首次将推理能力融入到了医学大模型。

医学问答模型虽然已经从BioBERT、PubMedBERT一类的领域预训练模型发展到Med-PaLM、HuatuoGPT、DISC-MedLLM、MedFound等生成式医学模型，但在真实应用场景中仍存在明显不足。其一，很多模型在考试题或标准问答集上表现较好，但面对日常慢病咨询中的口语化表达、连续追问和个性化问题时，稳定性仍然不足。其二，纯生成式模型往往难以同时满足“回答自然”和“证据清晰”这两个要求，尤其在治疗建议、药物说明和风险判断等问题上更容易暴露安全边界问题。其三，现有研究更多关注模型能力本身，而对知识更新、检索链路、证据展示和系统落地等系统实现环节关注不足。

\section{本文研究内容}

针对“三高”慢性病场景下大语言模型在可信性、安全性以及系统落地方面存在的问题\cite{JYRJ202407001}，本文围绕基于检索增强生成（Retrieval-Augmented Generation, RAG）的医学问答系统展开研究与实现。本文主要工作包括面向垂直医疗场景的关键技术设计以及系统工程实现。

本文的具体研究内容如下：

（1）围绕“三高”相关慢性病主题构建医学知识库，对相关政策文件、临床诊疗指南、健康科普资料以及互联网医学文本等数据进行整理和预处理，并通过文本切分和向量编码等方式，形成适用于检索增强生成框架的结构化知识基础。针对医学文本专业术语密集、信息表达高度压缩以及上下文依赖性较强等特点，重点分析了知识组织方式对后续检索效果与生成质量的影响。

（2）针对医学问答场景中单一检索策略难以同时兼顾精确匹配与语义召回的问题，设计了关键词检索与向量检索相结合的混合检索策略，并引入重排序与结果筛选机制，以提升检索证据的相关性与可靠性。在此基础上，构建基于大语言模型的生成流程，使模型能够结合检索到的外部医学证据进行答案生成，从而在一定程度上提升回答的准确性与可解释性。

（3）设计并实现了一套面向实际应用的医学问答系统，围绕用户问答交互、知识检索服务、模型推理调用以及结果展示等核心模块进行整体架构设计，实现了从知识导入、索引构建到在线问答服务的完整流程。同时，系统支持多轮对话、证据来源展示以及知识库管理等功能，使系统能够较好地服务于实际慢病健康咨询场景。

\section{论文组织结构}

本文共分为六章，各章内容安排如下：

第一章：绪论，介绍课题研究背景与意义，梳理国内外研究现状，并给出本文的研究内容与技术路线。

第二章：相关理论和技术介绍，主要阐述大语言模型、Transformer、检索增强生成（RAG）等基础理论，为后续方法设计提供理论支撑。

第三章：医学知识库构建与检索方法，重点介绍“三高”领域知识库的数据来源、预处理与向量化过程，以及BM25、向量检索与混合检索策略的设计与实验分析。

第四章：医学问答模型构建与优化，围绕RAG问答生成流程、Prompt与上下文构建、DPO偏好数据构造和训练过程展开，并通过实验评估模型优化效果。

第五章：系统设计与实现，从系统需求、总体架构、业务模块、数据库设计、模型推理服务与前端页面等方面，说明医学问答系统的工程实现过程。

第六章：：总结与展望，对全文工作进行归纳，分析当前工作的不足，并给出后续可改进方向。
